\textbf{Άσκηση 20}
\begin{thema}{Β}
Δίνεται η συνάρτηση $f(x) = \dfrac{e^x}{e^x-1}$.

\begin{erwthma}
\item Για να ορίζεται η συνάρτηση $f$, πρέπει ο παρονομαστής να είναι διάφορος του μηδενός:
\[ e^x - 1 \neq 0 \Leftrightarrow e^x \neq 1 \Leftrightarrow x \neq 0 \]
Συνεπώς, το πεδίο ορισμού της $f$ είναι το $D_f = \mathbb{R}^* = (-\infty, 0) \cup (0, +\infty)$. Υπολογίζουμε τα ζητούμενα όρια:
\begin{itemize}
    \item Για το $\lim\limits_{x \to 0^+} f(x)$:
    \[\lim\limits_{x \to 0^+} f(x) = \lim\limits_{x \to 0^+} \frac{e^x}{e^x-1} = \frac{1}{0^+} = +\infty\]
    \item Για το $\lim\limits_{x \to +\infty} f(x)$:
    \[ \lim_{x \to +\infty} \frac{e^x}{e^x-1} = \lim_{x \to +\infty} \frac{e^x}{e^x(1-e^{-x})} = \lim_{x \to +\infty} \frac{1}{1-e^{-x}} = \frac{1}{1-0} = 1 \]
\end{itemize}

\item Για να δείξουμε ότι η $f$ αντιστρέφεται, αρκεί να δείξουμε ότι είναι συνάρτηση $1-1$. Η $f$ είναι παραγωγίσιμη στο $D_f$ ως πηλίκο παραγωγίσιμων συναρτήσεων με:
\begin{align*}
f'(x) &= \left( \frac{e^x}{e^x-1} \right)' = \frac{(e^x)'(e^x-1) - e^x(e^x-1)'}{(e^x-1)^2} =\\&= \frac{e^x(e^x-1) - e^x \cdot e^x}{(e^x-1)^2} = \frac{e^{2x} - e^x - e^{2x}}{(e^x-1)^2} = \frac{-e^x}{(e^x-1)^2}<0
\end{align*}
για κάθε $x \in (-\infty, 0) \cup (0, +\infty)$.
\begin{itemize}
    \item Στο διάστημα $\Delta_1 = (-\infty, 0)$ η $f$ είναι γνησίως φθίνουσα, οπότε $f(\Delta_1) = \left(\lim\limits_{x \to 0^-} f(x), \lim\limits_{x \to -\infty} f(x)\right)$. Είναι 
\begin{itemize}[label=\tiny\faPlay]
\item $\lim\limits_{x \to -\infty} f(x) = \frac{0}{0-1} = 0$
\item $\lim\limits_{x \to 0^-} f(x) = \frac{1}{0^-} = -\infty$. Άρα $f(\Delta_1) = (-\infty, 0)$.
\end{itemize}
\item Στο διάστημα $\Delta_2 = (0, +\infty)$ η $f$ είναι γνησίως φθίνουσα, οπότε $f(\Delta_2) = \left(\lim\limits_{x \to +\infty} f(x), \lim\limits_{x \to 0^+} f(x)\right)$. Όπως έχουμε, $f(\Delta_2) = (1, +\infty)$.
\end{itemize}
Το σύνολο τιμών της $f$ είναι το $f(D_f) = f(\Delta_1) \cup f(\Delta_2) = (-\infty, 0) \cup (1, +\infty)$.
Επειδή η $f$ είναι γνησίως φθίνουσα σε καθένα από τα διαστήματα του πεδίου ορισμού της και τα σύνολα τιμών $f(\Delta_1), f(\Delta_2)$ είναι ξένα μεταξύ τους, η $f$ είναι $1-1$ στο $D_f$, άρα αντιστρέφεται.

Το πεδίο ορισμού της $f^{-1}$ είναι το σύνολο τιμών της $f$, δηλαδή $D_{f^{-1}} = (-\infty, 0) \cup (1, +\infty)$.
Για τον τύπο της $f^{-1}$, λύνουμε την εξίσωση $y = f(x)$ ως προς $x$:
\begin{align*}
y = \frac{e^x}{e^x-1} &\Leftrightarrow y(e^x-1) = e^x \Leftrightarrow y e^x - y = e^x \Leftrightarrow \\&\Leftrightarrow y e^x - e^x = y \Leftrightarrow e^x(y-1) = y \Leftrightarrow e^x = \frac{y}{y-1}\\&\Leftrightarrow
x = \ln \left( \frac{y}{y-1} \right)
\end{align*}
οπότε:
\[ f^{-1}(x) = \ln \left( \frac{x}{x-1} \right), \quad x \in (-\infty, 0) \cup (1, +\infty) \]

\item Οι ασύμπτωτες της $C_{f^{-1}}$ αναζητούνται στα άκρα του πεδίου ορισμού της $D_{f^{-1}} = (-\infty, 0) \cup (1, +\infty)$.
\begin{itemize}
    \item Στο $x=0$: $\lim\limits_{x \to 0^-} f^{-1}(x) = \lim\limits_{x \to 0^-} \ln \left( \frac{x}{x-1} \right)$.
    Θέτουμε $u = \frac{x}{x-1}$. Όταν $x \to 0^-$, τότε $u \to \frac{0}{-1} = 0^+$.
    Άρα $\lim\limits_{x \to 0^-} f^{-1}(x) = \lim\limits_{u \to 0^+} \ln u = -\infty$.
    Συνεπώς, η ευθεία $x=0$ είναι κατακόρυφη ασύμπτωτη της $C_{f^{-1}}$.
    \item Στο $x=1$: $\lim\limits_{x \to 1^+} f^{-1}(x) = \lim\limits_{x \to 1^+} \ln \left( \frac{x}{x-1} \right)$.
    Όταν $x \to 1^+$, τότε $x-1 \to 0^+$, οπότε $\frac{x}{x-1} \to +\infty$.
    Άρα $\lim\limits_{x \to 1^+} f^{-1}(x) = +\infty$.
    Συνεπώς, η ευθεία $x=1$ είναι κατακόρυφη ασύμπτωτη της $C_{f^{-1}}$.
    \item Στα $\pm\infty$: $\lim\limits_{x \to \pm\infty} f^{-1}(x) = \lim\limits_{x \to \pm\infty} \ln \left( \frac{x}{x-1} \right)$.
    Επειδή $\lim\limits_{x \to \pm\infty} \frac{x}{x-1} = 1$, το όριο είναι $\ln 1 = 0$.
    Συνεπώς, η ευθεία $y=0$ είναι οριζόντια ασύμπτωτη της $C_{f^{-1}}$ τόσο στο $+\infty$ όσο και στο $-\infty$.
\end{itemize}

\item Για $x > 0$, η $f$ έχει σύνολο τιμών $(1, +\infty)$, άρα $f(x) > 0$.
Το ζητούμενο εμβαδόν είναι:
\[ E = \int_{\ln 2}^{\ln 3} |f(x)| dx = \int_{\ln 2}^{\ln 3} \frac{e^x}{e^x-1} dx \]
Παρατηρούμε ότι ο αριθμητής είναι η παράγωγος του παρονομαστή, δηλαδή $e^x = (e^x-1)'$.
Επομένως:
\[ E = \int_{\ln 2}^{\ln 3} \frac{(e^x-1)'}{e^x-1} dx = \left[ \ln|e^x-1| \right]_{\ln 2}^{\ln 3} = \ln|e^{\ln 3}-1| - \ln|e^{\ln 2}-1| \]
\[ E = \ln|3-1| - \ln|2-1| = \ln 2 - \ln 1 = \ln 2 \text{ τ.μ.} \]
\end{erwthma}
\end{thema}
